% =====================================================================
% Section IX -- Discussion and Conclusion
% SpecAsync-UVM manuscript, draft 1 (2026-11-06)
% No new numbers: every number below appears in Sections I, VI-VIII and
% was checked there (C14/C15). Run the C15 number check over this file.
% [BUG] The stock-driver oops is described, not cited; if the NVIDIA issue
%       is filed before submission, cite its number here.
% =====================================================================

\section{Discussion and Conclusion}
\label{sec:discussion}

\subsection{What Off-Path Staging Can Do in This Driver}

Our results answer the question we set out to test, but not in the way the
idea's proponents, including our own earlier draft, expected. Moving residency
work off the fault-servicing thread works as an engineering matter: the worker
staged 99.95\% of pages ahead of demand, and the copies it made shortened the
lock-held dispatch. It does not remove faults, because only fault service
installs mappings (Section~\ref{sec:bg-resident-mapped}). Its one route to
preventing faults runs through the prefetcher's density count, and the same
effect is available, and larger, by lowering the threshold
(Section~\ref{sec:spec-matched}).

Two ways forward follow. The first is to stage from where mapping is possible.
Explicit prefetch from user space runs the driver's make-resident-and-map path,
which is why a user-space predictor can remove faults where our worker cannot
\cite{shen2026stride}. The second is to change the driver so that a background
path may install mappings. That is a larger change than ours: mapping touches
page tables the servicing thread also writes, and would have to respect the
driver's lock order and replay protocol. We have not attempted it. Our
measurements bound what it could gain only indirectly, since Stencil-24K
already runs 9.57\% faster when the prefetcher simply maps more on each fault.

\subsection{What a Better Threshold Would Need}

No fixed threshold dominates the shipped one (Section~\ref{sec:threshold}). Go
\emph{et al.}\ drew the same conclusion and built a controller that adapts the
threshold at run time \cite{go2023earlyadaptor}. Our measurements suggest three
inputs such a controller needs. The first is the access density within a
region: on oversubscribed managed memory and on HMM, the sign of the effect changed
between 8 and 64 pages per 2\,MB block. The second is memory pressure: on managed memory
the cost of a low threshold for sparse access was small in memory and large
once the array was oversubscribed. The third is the cost of migration itself:
on system-allocated memory, which moves in 4\,KB units, the same sparse access
was much slower at a low threshold even in memory. A controller tuned on
managed memory would therefore choose wrongly for HMM. These are observations
from one synthetic benchmark at one oversubscription ratio, and they locate a
crossover only between the densities we measured.

Part of the dense-access gain lies outside the GPU. The same threshold governs
prefetching for CPU faults on managed memory, and at threshold 0 the CPU faulted
six times less often while initialising the array; on dense cells 17--23\% of
the wall-clock gain was outside the kernels. A benchmark that initialises its
data on the GPU, or with explicit copies, would see less of it.

\subsection{System-Allocated Memory and Other Vendors}

HMM lets a GPU fault on ordinary memory, and Allen \emph{et al.}\ name it as a
motivation for studying the fault path \cite{allen2024finegrain}. On NVIDIA's
open driver it shares the managed path's prefetcher and threshold but not its
migration granularity, and that difference changed the sparse-access cost by
about two hundred percentage points (Section~\ref{sec:thr-hmm}). AMD GPUs reach system
memory through the same kernel interface, and Cooper \emph{et al.}\ report the
same tension between aggressive range prefetching and access pattern on
MI250X \cite{cooper2024svm}. Whether a density-aware prefetch policy would help
there, where the driver and its prefetcher are different, is a question our
measurements raise but cannot answer.

\subsection{Threats to Validity}

\paragraph{Scope of the speculation results}
The oracle is near-perfect only for Stencil-24K, whose first-touch order varies
by about a hundred ranks between runs; for GraphBFS-23 it is a weakened oracle.
The speculation results therefore bound what off-path staging can achieve on
one regular workload and do not evaluate a practical predictor. They come from
one GPU and one driver release.

\paragraph{Scope of the threshold results}
The dense-access threshold result replicates on a second GPU generation with
the same driver release. The sparse-access, oversubscription and HMM results
come from the RTX~5070~Ti alone, from one synthetic benchmark with one fixed
page list, and from one oversubscription ratio. The HMM results are for arrays
that fit in GPU memory: the first oversubscribed HMM run crashed the stock
driver, a null-pointer dereference in its chunk-free path followed by a soft
lockup of the fault handler, with no modified code loaded. The stop rule
forbade a retry, so the case was not measured.

\paragraph{Measurement}
The outcome of record is process wall-clock, which includes CUDA
initialisation and host-side setup; we chose it because it is what a user
experiences, and we separate kernel time from the rest where it matters
(Section~\ref{sec:thr-dense}). Each cell has ten runs, and comparisons across
sessions are reported side by side but never pooled. Two of our instruments
had limits we found only after use, GPU memory in use on managed memory and
the oracle's own lookup cost; both are in the catalog
(Section~\ref{sec:pitfalls}), and no claim rests on either.

\paragraph{The study's own history}
This paper reverses an earlier version of itself. The pre-registered gates,
the positive controls and the catalog are our response, and
Section~\ref{sec:pitfalls} reports the errors they caught and those they did
not.

\subsection{Conclusion}

We built off-path page staging into NVIDIA's open-source UVM driver and
measured it, with an oracle, under pre-registered tests. Staging cannot prevent
a fault, because in this driver only fault service maps a page. Against the
shipped driver it helps, by 3.4--5.3\% on a regular workload, only by raising
the count in the prefetcher's density rule, and lowering that rule's threshold
does more without any speculation. The threshold itself cuts both ways: it
speeds up dense access on two GPU generations and slows sparse access, under
memory pressure on managed memory and even without it on system-allocated
memory. For anyone building on the UVM fault path, the practical lesson is that
the lever worth adapting is the one the driver already has, and that the input
it most needs is how densely each region is used.
